% ================================================================
% chapters/ch1_introduction.tex
% Chapter 1 — Introduction
% ================================================================

\chapter{Introduction}
\label{ch:introduction}

% ================================================================
\section{Background}
\label{sec:background}
% ================================================================

Money has existed for thousands of years, but its form has
changed dramatically over time. From barter to coins, from
coins to paper notes, and now from paper notes to digital
payments --- each transition has brought new conveniences
alongside new challenges. The latest and perhaps most
significant transition is the move toward fully digital
currency issued directly by central banks, known as Central
Bank Digital Currency, or CBDC.

A CBDC is a digital form of a country's official currency,
issued and backed by the central bank, just as physical
banknotes are. Unlike private digital payment platforms such
as eSewa or Khalti, a CBDC is legal tender --- it carries the
full authority of the government and the central bank behind
it. Unlike cryptocurrencies such as Bitcoin or Ethereum, a
CBDC is not decentralised. It is issued, managed, and
overseen by the central bank of the issuing country.

To understand why CBDCs are being explored worldwide, it helps
to understand the limitations of existing payment systems.
Physical cash is anonymous --- when a person pays for groceries
with a banknote, no record is created linking that person to
that purchase. However, cash is inefficient for large
transactions, easily lost or stolen, and impossible to use
for remote or online payments. Digital payments solve these
problems but introduce a new one: every digital transaction
creates a data record. This record includes who paid, how
much, when, where, and to whom. These details are collectively
referred to as transaction metadata.

Transaction metadata is enormously valuable --- to banks for
fraud detection, to governments for tax compliance, and to
regulators for anti-money laundering monitoring. However, the
same metadata that serves legitimate purposes also creates
serious privacy risks. If transaction metadata can be accessed
by unauthorised parties, or if it can be analysed to identify
individuals who were supposed to be anonymous, then digital
currency systems fundamentally undermine the privacy that
physical cash provided.

This is the core tension at the heart of every CBDC design:
how do you build a digital currency that is efficient,
compliant with regulations, and inclusive --- while also
protecting the privacy of the citizens who use it? There is
no simple answer, and different cryptographic approaches offer
different tradeoffs between privacy and transparency. The
choice of cryptographic scheme is therefore one of the most
consequential design decisions in any CBDC deployment.

Globally, over 130 countries are exploring or piloting CBDCs
as of 2024. China's digital yuan, known as the e-CNY, has
been piloted across multiple major cities and has processed
billions of yuan in transactions. The European Central Bank
is advancing the digital euro project, with privacy described
as a core design requirement. The Eastern Caribbean Currency
Union launched the DCash system, and the Bahamas operates the
Sand Dollar --- the world's first fully deployed retail CBDC.
These projects represent a broad consensus that digital
currency is not a distant possibility but an imminent reality
for financial systems worldwide.

Nepal is not isolated from this global movement. The Nepal
Rastra Bank (NRB), Nepal's central bank, has been examining
the feasibility of a domestic CBDC as part of its digital
financial infrastructure strategy. Any CBDC deployed in Nepal
would sit on top of existing payment infrastructure, interact
with platforms like eSewa and Khalti that millions of Nepalis
already use, and serve a population that spans geographically
diverse zones from the dense urban concentration of the
Kathmandu valley to the remote hill and Terai rural regions.
The privacy implications of a CBDC in this specific context
have not yet been formally studied, and that is precisely the
gap this thesis addresses.

% ================================================================
\section{Nepal Context}
\label{sec:nepal_context}
% ================================================================

Nepal presents a distinctive and important context for CBDC
deployment that differs significantly from both Western
economies and the larger Asian economies typically studied in
the CBDC literature. Understanding Nepal's specific
characteristics is essential for understanding why this
thesis approaches the privacy problem the way it does.

\subsection{Financial Inclusion and the Unbanked Population}

Despite significant progress in recent years, a substantial
portion of Nepal's population remains either unbanked or
underbanked. Access to formal financial services is
geographically uneven, with urban centres having far greater
branch density than rural hill and Terai regions. A CBDC
represents a significant opportunity to extend financial
access to populations who have historically been excluded
from the formal banking system, particularly if it is
delivered through mobile platforms that do not require a
physical branch visit.

However, financial inclusion and privacy are not automatically
aligned. The very act of bringing previously unbanked citizens
into a digital payment system creates a new metadata record
for every transaction they make. For a person who has never
been part of a formal financial system, this represents a
sudden and significant increase in their digital footprint.
If the CBDC system does not adequately protect metadata
privacy, the populations that benefit most from financial
inclusion may also bear the greatest privacy risk.

\subsection{Mobile-First Payment Infrastructure}

Nepal's digital payment ecosystem has developed rapidly and
is predominantly mobile-first. eSewa and Khalti, Nepal's two
dominant digital wallet platforms, together process the
majority of digital transactions in the country. QR code
payments have been adopted widely at merchants ranging from
large supermarkets to small roadside stalls. Near-field
communication (NFC) payments are emerging in urban areas,
while offline token systems are being explored for low
connectivity rural deployment.

This mobile-first reality means that in a CBDC deployment,
the payment channel --- whether a transaction was made via
mobile wallet, QR code, NFC, or offline token --- is a
metadata field that is observable at the network level. As
this thesis will demonstrate, this attribute is relevant to
privacy risk modelling in Nepal's specific context and has
been incorporated into the model accordingly.

\subsection{Geographic Stratification}

Nepal's geography creates one of the most stratified
transaction environments of any CBDC deployment context
being considered globally. The country spans five distinct
geographic zones with dramatically different population
densities, merchant concentrations, and transaction
frequencies:

\begin{itemize}
    \item \textbf{Kathmandu Valley} --- the capital region,
    accounting for approximately 35\% of all digital
    transactions despite being a small geographic area.
    High merchant density, high transaction frequency.

    \item \textbf{Hill Urban} --- secondary cities such as
    Pokhara, Biratnagar, and Dharan. Moderate transaction
    density and growing digital payment adoption.

    \item \textbf{Hill Rural} --- the largest rural segment
    by population. Low merchant density, low transaction
    frequency, and limited connectivity in many areas.

    \item \textbf{Terai Urban} --- urban centres in Nepal's
    southern plains. Growing commercial activity and
    cross-border trade influence.

    \item \textbf{Terai Rural} --- rural southern plains
    with significant agricultural activity and informal
    economic transactions.
\end{itemize}

This geographic stratification matters for privacy because
transaction zone is an observable metadata field in any
CBDC system with location-aware merchant registration.
A transaction recorded in a hill rural zone with limited
merchants may be far more uniquely identifying than an
identical transaction recorded in the Kathmandu valley,
simply because there are fewer possible users that
transaction could belong to. This thesis incorporates
transaction zone as a Nepal-specific attribute in its
model for precisely this reason.

\subsection{Constitutional Right to Privacy}

Nepal's Constitution, promulgated in 2015, guarantees the
right to privacy as a fundamental right under Article 28.
This constitutional protection extends to personal
information and communications, and creates a legal
obligation for any state-sponsored digital system to
protect citizen privacy. A CBDC issued by the Nepal Rastra
Bank would be a state-sponsored system, and its privacy
properties are therefore not merely a technical preference
but a constitutional requirement.

\subsection{The Nepal Rastra Bank's Digital Currency Exploration}

The Nepal Rastra Bank has been actively studying digital
financial infrastructure, including CBDC feasibility, as
part of its broader financial modernisation agenda. NRB
has noted the importance of privacy, security, and
financial inclusion as core requirements for any digital
currency initiative. The deployment platform most
frequently referenced in the context of permissioned
blockchain-based CBDC is Hyperledger Fabric, a
production-grade enterprise blockchain framework that
provides the permissioned identity management and
chaincode-based transaction processing capabilities
required for a regulated financial system.

This thesis models the privacy risk specifically within
the context of a Hyperledger Fabric-based token CBDC
under NRB's domestic jurisdiction. Cross-border
remittances, while significant to Nepal's economy,
involve separate settlement rails and regulatory
frameworks and are therefore outside the scope of this
domestic model.

% ================================================================
\section{The Problem}
\label{sec:problem}
% ================================================================

The problem this thesis addresses can be stated plainly:
when a CBDC transaction is recorded on a digital ledger,
it leaves behind a set of observable metadata fields.
These fields --- the amount, the time, the merchant, the
wallet identifier, the transaction frequency, the payment
channel, and the geographic zone --- do not contain the
user's name. But the combination of these fields may be
sufficient to identify the user anyway.

This is not a hypothetical concern. In 2006, researchers
Narayanan and Shmatikoff demonstrated that supposedly
anonymous Netflix viewing records could be linked to real
individuals by correlating them with public ratings on
the Internet Movie Database. The users had no names
attached to their Netflix records. Yet their viewing
patterns alone were unique enough that knowing just eight
ratings was sufficient to identify 99\% of users in the
dataset. The lesson was direct: removing a name from a
record does not make the record anonymous if the
remaining data is sufficiently unique.

The same principle applies to CBDC transaction metadata.
A user who pays NPR 4,738.50 at a specific merchant at
2:47am using a mobile wallet in a hill rural zone has
created a transaction fingerprint that may be unique to
exactly one person in the entire population. No
cryptographic attack is needed. No key is broken. The
information is simply there in the metadata, waiting to
be read.

\subsection{The Pseudonymity Problem}

Hyperledger Fabric, the platform most relevant to Nepal's
CBDC context, operates with a permissioned identity model.
Every participant has an X.509 certificate issued by the
Membership Service Provider (MSP). The Nepal Rastra Bank,
as the network operator, knows the real-world identity
behind every wallet address. This is by design --- KYC
and AML compliance requires it.

However, this does not mean the privacy problem disappears.
The ledger is visible to all peers on a channel, which
may include financial institutions, payment processors,
and merchant organisations who are channel participants
but who should not have access to the complete transaction
history of every citizen. Furthermore, if the transaction
log is ever separated from the identity mapping --- through
a data breach, an insider threat, or a regulatory
disclosure --- the richness of the metadata alone may be
sufficient to re-identify users without ever consulting
the identity table.

This is the pseudonymity problem. A CBDC system may
replace real names with wallet addresses, but if the
transaction metadata is rich enough, the pseudonym itself
becomes identifying. The wallet address of a user who
always pays at the same four merchants, always at the
same times, and always in amounts consistent with a
specific income level is effectively a name --- it just
uses numbers instead of letters.

\subsection{The Measurement Gap}

Despite the significance of this problem, no formal
quantitative framework exists for measuring metadata
de-anonymization risk across different CBDC cryptographic
scheme choices. Existing comparisons between schemes are
qualitative: researchers and policymakers say that ring
signatures provide more privacy than blind signatures,
or that zero-knowledge proofs provide stronger guarantees
than hash-based tokens. These statements are directionally
correct but provide no quantitative basis for decision-making.

A policymaker at NRB considering whether to deploy Ring
Signature tokens or ZKP tokens cannot currently answer
questions such as:

\begin{itemize}
    \item How much more privacy does Ring Signature
    provide compared to Blind Signature, expressed as
    a measurable quantity?

    \item If Ring Signature cannot be deployed due to
    computational constraints, how much privacy is lost
    by falling back to Blind Signature, and can that
    loss be recovered through mitigation strategies?

    \item At Nepal's target deployment scale of
    100,000 users, does the risk profile of each scheme
    change significantly compared to a pilot of 1,000
    users?

    \item Are the results of any such analysis stable,
    or do they depend on the specific data used to
    produce them?
\end{itemize}

Without answers to these questions, scheme selection
is either arbitrary or driven by factors other than
privacy --- such as implementation cost or familiarity
--- with no formal basis for claiming that the chosen
scheme provides adequate privacy protection.

This thesis provides the mathematical framework,
the implementation, and the empirical analysis required
to answer all of these questions with specific numbers.

% ================================================================
\section{Research Question}
\label{sec:research_question}
% ================================================================

The central research question motivating this thesis is:

\vspace{0.5em}
\begin{center}
\textit{How can the metadata de-anonymization risk of
different cryptographic schemes in a token-based CBDC
be formally measured, normalised, and compared on a
common quantitative scale?}
\end{center}
\vspace{0.5em}

This primary question gives rise to four secondary
research questions that structure the analysis:

\begin{enumerate}

    \item \textbf{Scheme comparison.} Across six
    cryptographic schemes --- Plain Hash Token, Blind
    Signature, Hybrid Scheme, Ring Signature, Stealth
    Address, and ZKP Token --- what are the relative
    metadata privacy risk levels at a baseline population
    of $N = 1{,}000$ users?

    \item \textbf{Population scalability.} How does the
    risk profile of each scheme change as the CBDC
    population grows from a small village pilot of
    $N = 10$ users to a national deployment of
    $N = 100{,}000$ users? Does scale improve privacy,
    and if so, by how much and for which schemes?

    \item \textbf{Mitigation effectiveness.} How
    effectively do privacy-enhancing technical
    mitigations --- differential privacy on amounts,
    timestamp suppression, wallet unlinkability, and
    Nepal-specific controls --- reduce the risk score
    of each scheme, and which mitigations provide the
    greatest return?

    \item \textbf{Nepal-specific context.} How do
    Nepal-specific metadata attributes --- payment
    channel and transaction zone --- affect the risk
    model, and what does the model reveal about the
    privacy implications of Nepal's specific demographic
    and geographic structure?

\end{enumerate}

These questions are answered through the development
and application of a formal mathematical model whose
design, derivation, and implementation are described
in detail in Chapter~\ref{ch:methodology}, and whose
results are presented and discussed in
Chapter~\ref{ch:results}.

% ================================================================
\section{Research Objectives}
\label{sec:objectives}
% ================================================================

The overarching objective of this thesis is to develop a
formal, quantitative framework for measuring metadata
de-anonymization risk in token-based CBDC systems and to
apply that framework in Nepal's specific deployment context.
This overarching objective is broken down into five concrete
objectives that together define the scope and contribution
of the work.

The first objective is to derive a normalised privacy risk
metric from mathematical first principles. This metric,
denoted $R_{\mathrm{norm}}$, must be bounded in the interval
$[0, 1]$, must be comparable across different numbers of
attributes and different population sizes, and must have a
proven theoretical upper bound that gives the normalisation
its meaning. Achieving this requires deriving both the
raw risk accumulation formula and the theoretical maximum
against which it is normalised.

The second objective is to model six cryptographic schemes
that represent the realistic spectrum of token-based CBDC
privacy designs. These schemes --- Plain Hash Token, Blind
Signature, Hybrid Scheme, Ring Signature, Stealth Address,
and ZKP Token --- span from the weakest to the strongest
available privacy guarantees. Each scheme must be modelled
in terms of how it transforms the value domain of each
observable metadata attribute, so that differences in
$R_{\mathrm{norm}}$ across schemes reflect genuine
cryptographic design differences rather than arbitrary
parameter choices.

The third objective is to incorporate Nepal-specific context
into the model. This means identifying the metadata
attributes most relevant to Nepal's domestic CBDC deployment,
grounding their value distributions in Nepal's payment
infrastructure and geographic structure, and formally
justifying the exclusion of attributes outside the domestic
scope.

The fourth objective is to conduct a comprehensive empirical
analysis covering eight distinct experimental conditions:
baseline scheme comparison, population scalability from
village pilot to national scale, mitigation strategy
effectiveness, attribute sensitivity as the number of modelled
attributes grows, statistical robustness across multiple
random seeds, micro-scale behaviour at very small populations,
sensitivity of results to the interaction coefficient, and
worst-case scenario calibration against the theoretical
maximum.

The fifth objective is to translate the quantitative findings
into actionable policy recommendations for the Nepal Rastra
Bank, providing a formal basis for cryptographic scheme
selection that can inform NRB's CBDC deployment decision.

% ================================================================
\section{Scope and Limitations}
\label{sec:scope}
% ================================================================

Every research model operates within a defined boundary.
Being explicit about that boundary is not a weakness --- it
is a requirement for honest and reproducible science. This
section defines what this thesis does and does not cover,
and why each boundary was drawn where it was.

\subsection{Scope}

This thesis models a domestic, token-based CBDC under the
jurisdiction of the Nepal Rastra Bank. The reference
deployment platform is Hyperledger Fabric, a permissioned
enterprise blockchain framework. The model operates on
transaction metadata --- the observable fields associated
with a completed transaction --- and measures the risk that
this metadata can be used to re-identify users.

The model covers seven metadata attributes: transaction
amount, timestamp, merchant identifier, wallet address,
transaction frequency, payment channel, and transaction
zone. These seven attributes correspond to fields that would
be observable on a Hyperledger Fabric ledger in a realistic
CBDC deployment. The model covers six cryptographic schemes,
five mitigation strategies, populations ranging from
$N = 10$ to $N = 100{,}000$, and attribute sets ranging
from $m = 1$ to $m = 7$.

\subsection{Limitations}

The following limitations are acknowledged and should be
considered when interpreting results.

\textbf{Synthetic data.} The model uses synthetically
generated transaction data rather than real NRB transaction
records. The distributions chosen for each attribute are
informed by Nepal's payment infrastructure context, but
they are approximations. Results may differ if applied to
real transaction data with different statistical properties.

\textbf{Domestic scope only.} Cross-border remittance
transactions are excluded. While remittances constitute
a significant portion of Nepal's financial flows,
they involve separate settlement rails, foreign
correspondent banking relationships, and regulatory
frameworks outside NRB's domestic CBDC mandate.
Remittance corridor was considered as an attribute
and explicitly excluded on these grounds.

\textbf{Pairwise interactions only.} The model captures
interaction effects between pairs of attributes. Higher-order
interactions --- triplets, quadruplets, and beyond --- are
not modelled. For most practical deployment scenarios,
pairwise interactions capture the dominant amplification
effects. However, in datasets with strong higher-order
correlations, the model may underestimate total risk.

\textbf{Passive adversary only.} The threat model assumes
a passive observer who reads the ledger but does not
modify it, inject transactions, or perform active attacks
such as Sybil attacks or transaction injection. Active
attack vectors are outside the scope of this model.

\textbf{Hyperledger Fabric specific fields.} Several
Fabric-specific metadata fields are not included in the
model, including transaction ID linkability, endorsing
peer identity, chaincode version, and X.509 certificate
metadata. These fields represent additional leakage channels
that would increase $m$ and consequently $R_{\mathrm{max}}$
if included. Their exclusion is acknowledged as a limitation
and represents a direction for future work.

\textbf{Static population.} The model assumes a fixed
population of $N$ users. In practice, a CBDC user base
grows over time as adoption increases. The scalability
analysis in Chapter~\ref{ch:results} partially addresses
this by testing across a wide range of $N$ values, but
does not model dynamic user entry and exit.

% ================================================================
\section{Significance of the Study}
\label{sec:significance}
% ================================================================

The significance of this work operates at three levels:
academic, policy, and societal.

At the academic level, this thesis makes a contribution
to the formal privacy measurement literature by introducing
a normalised metric specifically designed for the CBDC
metadata context. Existing privacy frameworks --- k-anonymity,
l-diversity, and differential privacy --- were not developed
with CBDC metadata in mind and do not directly address the
multi-scheme comparison problem this thesis solves.
The derivation of $R_{\mathrm{max}}(m) = m + \binom{m}{2}$
as a proven theoretical upper bound, and the use of this
bound to produce a scale-independent normalised score,
represents a formal contribution that extends the existing
privacy measurement literature into a new application domain.

At the policy level, this thesis provides the Nepal Rastra
Bank with the first quantitative tool for comparing CBDC
cryptographic schemes on the basis of metadata privacy risk.
Prior to this work, any comparison between schemes relied
on qualitative judgement. This thesis replaces that
judgement with specific numbers: $R_{\mathrm{norm}} = 0.317$
for Plain Hash Token and $R_{\mathrm{norm}} = 0.035$ for
ZKP Token at $N = 1{,}000$, with the full risk ladder,
mitigation effectiveness data, and population scalability
curves available for any deployment scenario NRB wishes
to evaluate.

At the societal level, the stakes of getting CBDC privacy
right are real and affect ordinary people. Nepal has a
constitutional right to privacy, a large rural population
that will be entering the digital financial system for
the first time, and a geographic structure that makes
certain communities particularly vulnerable to
re-identification from transaction metadata. The
micro-scale analysis in this thesis demonstrates that
at a village pilot scale of $N = 50$ users, all schemes
except ZKP Token produce $R_{\mathrm{norm}} > 0.48$ ---
meaning nearly half of pilot participants are potentially
re-identifiable from their transaction metadata alone.
This finding has direct implications for how NRB should
approach any community-level CBDC pilot, regardless of
which scheme is ultimately selected for national deployment.

Furthermore, Nepal is not alone in this context. Many
developing economies share Nepal's characteristics ---
mobile-first infrastructure, geographic stratification,
high remittance dependency, and a large unbanked population
entering the digital economy. The framework developed in
this thesis is generalizable beyond Nepal. Any central bank
facing similar deployment conditions can substitute their
own attribute set and population parameters into the model
and obtain equivalent analysis. This positions the work
as a contribution not only to Nepal's CBDC policy but
to the broader global conversation about how to build
digital currencies that serve financial inclusion without
sacrificing the privacy of the people they are designed
to help.

% ================================================================
\section{Report Structure}
\label{sec:structure}
% ================================================================

This report is organised into five chapters, each building
on the previous to construct a complete argument from
problem identification through to policy recommendation.

Chapter~\ref{ch:introduction} --- the present chapter ---
establishes the background, context, and motivation for
the research. It introduces the CBDC privacy problem,
situates it within Nepal's specific deployment context,
defines the research questions and objectives, and
honestly states the scope and limitations of the work.

Chapter~\ref{ch:literature} reviews the existing body of
knowledge relevant to this thesis. It covers the global
CBDC landscape, the concept of privacy in digital payment
systems, the cryptographic schemes evaluated in this work,
known de-anonymization attacks including the foundational
Netflix attack, existing formal privacy measurement
frameworks and their limitations, and the role of
Hyperledger Fabric as the reference platform. The chapter
concludes by identifying the specific research gap that
this thesis fills.

Chapter~\ref{ch:methodology} presents the complete
theoretical framework and implementation. It defines the
system model and adversary model, derives the mathematical
components of the risk metric from first principles,
describes the six cryptographic schemes and their metadata
transformations, justifies the Nepal-specific attribute
set, and documents the synthetic data generation approach
and Python implementation.

Chapter~\ref{ch:results} presents and discusses the
results of eight analyses. Each analysis is presented
with its motivating question, methodology, results, and
interpretation. The chapter covers baseline scheme
comparison, population scalability, mitigation
effectiveness, attribute sensitivity, seed sensitivity,
micro-scale pilot behaviour, lambda sensitivity, and
worst-case calibration.

Chapter~\ref{ch:conclusions} synthesises the findings,
provides direct answers to the research questions stated
in this chapter, summarises the contributions of the
thesis, makes specific recommendations to the Nepal
Rastra Bank, and identifies directions for future work
that extend or build upon this model.